\chapter{Competition creates competition}\label{ch:ccc}

\begin{chapterplan}
Why can a stronger incumbent raise the probability that a challenger prepares, when it never raises what
preparation earns at any belief? This chapter states Proposition~2, reads its three conditions, and rebuilds the
proof economy by economy: orders, then posteriors, then entry, then prices.
\end{chapterplan}

\section{Three economies in words}\label{sec:ccc-story}

At fixed information, a stronger incumbent hurts a challenger: it never raises \(g_H\) or \(g_L\)
(Proposition~1(a)), so gross profit \(B(\mu)\) never rises at any fixed belief. But it also never narrows the
target-payoff spread \(\DT=t_H-t_L\) (Proposition~1(b)), on which the informed investor trades. The result
compares three economies that differ only in the incumbent: a weak \(F_0\), a strong \(F_1\) that strengthens
\(F_0\), and a very strong \(F_2\) that strengthens \(F_1\).

\begin{intuition}[two forces with one cause]
Against \(F_0\) the spread is below the trading cost, the price is silent, and only a zero-cost challenger
prepares. Against \(F_1\) the investor trades, and a purchase lifts the belief from \(\pi\) to \(M(\pi)\), where
gross profit covers \(c_H\); this jump outweighs the fall in profit, and entry rises. Against \(F_2\) information
is already on, so only the fall in profit acts: no price covers \(c_H\), and entry returns to \(\rho\).
\end{intuition}

\section{The objects of the result}\label{sec:ccc-setup}

The proof uses these objects from Chapters~\ref{ch:auction} to~\ref{ch:trading}. Under the class ranking,
\(\DT\ge0\) and \(g_H\ge g_L\), so \(B(\mu)=g_L+\mu(g_H-g_L)\) never falls in \(\mu\). Entry is
\(e(x)=\rho+(1-\rho)\zeta(P(x))\ge\rho\); \(b=(1-\lambda)\beta\), \(b_0=(1-\lambda)(1-2\beta)\), \(u=\lambda/b\). After a
nonzero order the posterior lies in \([m(\pi),M(\pi)]\) for every \(\beta<1/2\) (Lemma~A.1(b)). The price formula
is \(P(x)=t_0+e(x)\,w(\mu(x))\) with \(w(\mu)=t_L-t_0+\DT\mu\ge0\). The trading returns
\(K_H=e(1)\DT[1-\mu(1)]\) and \(K_L=e(-1)\DT\mu(-1)\) obey, under every candidate,
\(\rho\kappa(\pi)\DT\le K_j\le\DT\); a wrong-signed order earns at most \(-k\). A \emph{candidate} meets
(E1)--(E4) and (E6); an equilibrium also meets (E5). Against \(F_i\) write \(\Delta_i\), \(B_i\), \(g_{j,i}\); in
proofs \(M=M(\pi)\), \(m=m(\pi)\), \(\kappa=\kappa(\pi)\).

\section{The statement}\label{sec:ccc-statement}

\begin{paperresult}{Proposition 2 (Competition creates competition)}
Let the classes satisfy the class ranking and have finite means, let \(F_1\) be a continuous first-order
stochastic strengthening of the continuous incumbent distribution \(F_0\), and let \(F_2\) strengthen \(F_1\) in
the same sense. Fix \(0<\pi,\rho,\lambda<1\), \(0<\beta<1/2\), \(c_H>0\) and \(k>0\). Write \(\Delta_i=\DT(F_i)\),
\(B_i=B_{F_i}\), and \(t_{0,i}\) and \(t_{L,i}\) for the proceeds \(t_0\) and \(t_L\) against \(F_i\), with
\(b=(1-\lambda)\beta\) and \(\kappa(\pi)=\min\{m(\pi),1-M(\pi)\}\). Suppose that
\[
\begin{array}{ll}
\text{(H1), the high-cost window:}& B_0(\pi)<c_H<B_1(M(\pi)),\\
\text{(H2), the trading-cost window:}& \Delta_0<k<\rho\kappa(\pi)\Delta_1,\\
\text{(H3), the exclusion condition:}& B_2(M(\pi))<c_H.
\end{array}
\]
Each economy has a unique equilibrium outcome under all four tie rules.
\begin{enumerate}[label=(\roman*),leftmargin=2em]
\item \emph{The weak incumbent \(F_0\).} Both classes abstain, every posterior equals \(\pi\), entry is \(\rho\),
  and all orders share the price \(t_{0,0}+\rho[t_{L,0}-t_{0,0}+\Delta_0\pi]\). High-class ownership is
  \(\pi\rho W_H(F_0)\).
\item \emph{The strong incumbent \(F_1\).} Class \(H\) buys and class \(L\) sells with probability one. The
  posteriors after \((-1,0,1)\) are \((m(\pi),\pi,M(\pi))\), and entry is \((\rho,\rho,1)\). Prices rise strictly in
  this order. Class entry and total entry are
  \[
  e_H=\rho+(1-\rho)(b+\lambda),\qquad e_L=\rho+(1-\rho)b,\qquad \Ent_1=\rho+(1-\rho)(b+\pi\lambda)>\rho.
  \]
  High-class ownership rises exactly when \(e_HW_H(F_1)>\rho W_H(F_0)\). Class \(H\)'s takeover share rises exactly
  when
  \begin{equation}\label{eq:ccc-share}
  e_HW_H(F_1)W_L(F_0)>e_LW_H(F_0)W_L(F_1).
  \end{equation}
  If class \(H\) outbids the stronger incumbent, \(W_H(F_1)=1\), it also outbids \(F_0\); high-class ownership rises
  from \(\pi\rho\) to \(\pi e_H\), and its takeover share is nondecreasing, strictly exactly when \(W_L(F_0)>0\). If
  \(W_L(F_0)=0\), that share is one against both incumbents.
\item \emph{The very strong incumbent \(F_2\).} The orders and posteriors are those of part~(ii), entry is \(\rho\)
  after every order, and the three prices differ. High-class ownership is \(\pi\rho W_H(F_2)\).
\end{enumerate}
Against each \(F_i\), with \(\mu_i(x)\) and \(e_i(x)\) the posterior and entry after the order \(x\), prices are
\(P_i(x)=t_{0,i}+e_i(x)[t_{L,i}-t_{0,i}+\Delta_i\mu_i(x)]\), the price belief is \(\mu_i(x)\), high-cost preparation
is \([e_i(x)-\rho]/(1-\rho)\), and bids are truthful.
\end{paperresult}

\section{The three conditions in words}\label{sec:ccc-conditions}

Each inequality has one job (Appendix A.3 of the paper). \(B_0(\pi)<c_H\) gives entry \(\rho\) against \(F_0\) and,
with Proposition~1(a), no recruitment after the zero order against \(F_1\). \(c_H<B_1(M(\pi))\) gives recruitment
after a purchase against \(F_1\). \(\Delta_0<k\) stops every informed order against \(F_0\). \(k<\rho\kappa(\pi)\Delta_1\)
forces full orders against \(F_1\) and, with Proposition~1(b), against \(F_2\). \(B_2(M(\pi))<c_H\) stops recruitment
against \(F_2\); without it, parts (i) and (ii) still hold, but entry need not fall. If \(\Delta_0>k\), the price can carry
information against \(F_0\), and Chapter~\ref{ch:coexistence} shows strengths with two equilibria.

\begin{watchout}
Proposition~2 allows \(0<\beta<1/2\), while the exclusion bound of Chapter~\ref{ch:beliefs} needs \(\beta\le1/3\)
to bound the posterior after the zero order under every strategy. The proof avoids that bound: it first fixes
full orders with the bounds after nonzero orders, so the posterior after the zero order is exactly \(\pi\).
\end{watchout}

\section{When the trading-cost window is nonempty}\label{sec:ccc-window}

Some \(k\) satisfies (H2) exactly when \(\Delta_0<\rho\kappa(\pi)\Delta_1\). Since \(0\le\Delta_0\le\Delta_1\), this
needs \(\Delta_1>0\), and it reads \(\Delta_0/\Delta_1<\rho\kappa(\pi)\): a condition with no unit of money. Next,
\(1-M(\pi)=m(1-\pi)\), and \(m\) rises strictly (its odds are \(\pi/[(1-\pi)(1+u)]\)), so
\[
\kappa(\pi)=m\bigl(\min\{\pi,1-\pi\}\bigr)\le m(1/2)=\frac{1}{2+u},
\]
with equality only at \(\pi=1/2\). A prior away from one half narrows the window: the price already leans toward
the favored class, so that class's order tells the market little. Since \(\rho\kappa(\pi)<1/2\), a nonempty window
needs \(\Delta_1>2\Delta_0\).

\section{The proof}\label{sec:ccc-proof}

\begin{keyidea}
In each economy the proof fixes the investor's orders first, with a bound that holds for every candidate. The
orders give the posteriors, the posteriors give entry, and entry and posteriors give the prices. Because the first
bound holds for every candidate, the tuple that comes out is the only equilibrium. A direct check of (E1)--(E6) at
the tuple then shows that it exists.
\end{keyidea}

\begin{derivation}[Step 0: inequalities that the conditions imply]
(a) Proposition~1(a) gives \(B_1\le B_0\) at every belief, so \(B_1(\pi)\le B_0(\pi)<c_H<B_1(M)\).
(b) The odds of \(M\) and \(m\) are \((1+u)\) times and \(1/(1+u)\) times the prior odds, so \(m<\pi<M\).
(c) Linearity gives \(0<B_1(M)-B_1(\pi)=(M-\pi)(g_{H,1}-g_{L,1})\), so \(g_{H,1}>g_{L,1}\), the threshold \(\tau_1\) exists, and \(\pi<\tau_1<M\).
(e) From (H2), \(\Delta_1>k/(\rho\kappa)>0\); Proposition~1(b) gives \(\rho\kappa\Delta_2\ge\rho\kappa\Delta_1>k\).
(f) For every \(\mu\le M\), \(B_2(\mu)=B_2(M)-(M-\mu)(g_{H,2}-g_{L,2})\le B_2(M)<c_H\) by (H3); this needs no
threshold \(\tau_2\), which may not exist.
(g) In every candidate, \(e(x)\ge\rho\).
\end{derivation}

\begin{derivation}[Step 1, part (i): orders, beliefs, entry and price against \(F_0\)]
Fix any candidate. The trading bound and the left side of (H2) give \(U_H(1)=K_H-k\le\Delta_0-k<0=U_H(0)\),
and in the same way \(U_L(-1)<0\); wrong-signed orders earn at most \(-k\). So both classes abstain,
\(\sigma_H=\sigma_L=\delta_0\), with no indifference. Then \(f_H(x)=f_L(x)\) for every order, so Bayes' rule gives
\(\mu(x)=\pi\) after every order, and the belief at every price is \(\pi\), whatever orders share it. Since
\(B_0(\pi)<c_H\) strictly, \(\zeta=0\) and \(e(x)=\rho\). The price formula gives one price
\(\bar P_0=t_{0,0}+\rho[t_{L,0}-t_{0,0}+\Delta_0\pi]\). Preparation is independent of \((\theta,R)\) given the
class, so \(\Own_{H,0}=\pi\rho\Pr(\theta>R\mid H)=\pi\rho\,\E_H[F_0(\theta)]=\pi\rho\,W_H(F_0)\).
\end{derivation}

\begin{derivation}[Steps 2 and 3, part (ii): orders, posteriors, entry, prices against \(F_1\)]
\emph{Orders.} Fix any candidate. By Step~0(g), Lemma~A.1(b) and \(\Delta_1>0\),
\begin{align*}
K_H&=e(1)\Delta_1[1-\mu(1)]\ge\rho\Delta_1(1-M)\ge\rho\kappa\Delta_1>k,\\
K_L&=e(-1)\Delta_1\mu(-1)\ge\rho\Delta_1m\ge\rho\kappa\Delta_1>k \why{right side of (H2)}.
\end{align*}
So \(U_H(1)>0=U_H(0)>-k\ge U_H(-1)\), and the mirror for class \(L\): \(\sigma_H=\delta_1\), \(\sigma_L=\delta_{-1}\),
with no indifference. This uses only nonzero orders, so it holds for every \(\beta<1/2\).

\emph{Posteriors.} On \((-1,0,1)\), \(f_H=(b,b_0,b+\lambda)\) and \(f_L=(b+\lambda,b_0,b)\). The likelihood ratios
are \(1/(1+u)\), \(1\), \(1+u\), so the posteriors are \((m,\pi,M)\) (Lemma~A.1(d)).

\emph{The price reveals the posterior.} If two orders shared one price, they would share entry \(\bar e\ge\rho>0\),
and \(0=P(x)-P(x')=\bar e\Delta_1[\mu(x)-\mu(x')]\) would force equal posteriors, which is false. So
\(\bar\mu(P(x))=\mu(x)\).

\emph{Entry.} \(B_1(M)>c_H\) gives \(e(1)=1\). \(B_1(\pi)<c_H\) gives \(e(0)=\rho\). And
\(B_1(m)=B_1(\pi)-(\pi-m)(g_{H,1}-g_{L,1})<c_H\) gives \(e(-1)=\rho\). All comparisons are strict.

\emph{Prices.} The increment identity (Appendix A.2 of the paper) is
\begin{align*}
P(x)-P(x')=e(x)\,\DT\,[\mu(x)-\mu(x')]+[e(x)-e(x')]\,w(\mu(x')).
\end{align*}
So \(P(0)-P(-1)=\rho\Delta_1(\pi-m)>0\) and \(P(1)-P(0)=\Delta_1(M-\pi)+(1-\rho)w_1(\pi)>0\), since \(w_1\ge0\).
\end{derivation}

\begin{derivation}[Step 4: entry, ownership and the takeover share]
Given the order, preparation does not depend on the class, and \(2b+b_0+\lambda=1\). So
\begin{align*}
e_H&=b\rho+b_0\rho+(b+\lambda)=\rho+(1-\rho)(b+\lambda),\qquad e_L=(b+\lambda)\rho+b_0\rho+b=\rho+(1-\rho)b,\\
\Ent_1&=\pi e_H+(1-\pi)e_L=\rho+(1-\rho)(b+\pi\lambda)>\rho \why{\(b>0\)}
\end{align*}
and \(e_H-e_L=(1-\rho)\lambda>0\). As in Step~1, \(\Own_{j,1}=\pi_je_jW_j(F_1)\) and \(\Own_{j,0}=\pi_j\rho W_j(F_0)\),
so ownership rises exactly when \(e_HW_H(F_1)>\rho W_H(F_0)\).

\emph{Shares.} Both denominators are positive: \(B_1(M)>c_H>0\) makes some \(g_{j,1}>0\), which forces
\(W_j(F_1)>0\), and \(W_j(F_0)\ge W_j(F_1)\). For shares, \(x/(x+y)>x'/(x'+y')\iff xy'>x'y\); with the ownership
formulas and division by \(\pi(1-\pi)\rho\), this is \eqref{eq:ccc-share}.

\emph{Outbidding.} If \(W_H(F_1)=1\), then \(1\ge W_H(F_0)\ge W_H(F_1)=1\), and \(\Own_{H,1}=\pi e_H>\pi\rho=\Own_{H,0}\).
The share condition becomes \(e_HW_L(F_0)>e_LW_L(F_1)\). If \(W_L(F_0)>0\),
\(e_HW_L(F_0)>e_LW_L(F_0)\ge e_LW_L(F_1)\), a strict rise. If \(W_L(F_0)=0\), then \(W_L(F_1)=0\) and the share is one
in both economies.

\emph{A sufficient condition.} The paper notes that ownership rises if
\((1+\rho)W_H(F_1)\ge2\rho W_H(F_0)>0\) and \(b+\lambda>\tfrac12\): then \(e_H>(1+\rho)/2\), so
\(e_HW_H(F_1)>\tfrac{1+\rho}{2}W_H(F_1)\ge\rho W_H(F_0)\).
\end{derivation}

\begin{watchout}
Do not state ``high-class ownership rises'' without a condition. A stronger incumbent can lower class \(H\)'s
chance to win. Ownership rises exactly when \(e_HW_H(F_1)>\rho W_H(F_0)\); the clean rise needs \(W_H(F_1)=1\).
\end{watchout}

\begin{derivation}[Step 5, part (iii): the very strong incumbent]
Fix any candidate. Step~0(e) gives \(\rho\kappa\Delta_2>k\), so the order argument of Step~2 gives full orders.
The order probabilities do not depend on the incumbent, so the posteriors are \((m,\pi,M)\), and every price belief
lies in \([m,M]\). Step~0(f) gives \(B_2<c_H\) there, so \(\zeta\equiv0\) and \(e\equiv\rho\). The increment
identity gives \(P(x)-P(x')=\rho\Delta_2[\mu(x)-\mu(x')]\), with \(\Delta_2\ge\Delta_1>0\), so the three prices
differ. As in Step~1, \(\Own_{H,2}=\pi\rho W_H(F_2)\).
\end{derivation}

\emph{Step 6: all four tie rules.} The investor rule matters only when \(K_j=k\); the challenger rule only when \(B(\bar\mu)=c_H\). Against each
\(F_i\), the comparison with \(k\) is strict for every candidate (Steps 1, 2 and 5). These comparisons fix the
orders in every equilibrium, so they fix the beliefs, and each comparison with \(c_H\) is strict too. No tie rule
is ever used, so all four give the same unique outcome. For the last paragraph of the statement, solve
\(e=\rho+(1-\rho)\zeta\) to get \(\zeta=(e-\rho)/(1-\rho)\).

\section{Decisive advantage}\label{sec:ccc-decisive}

\begin{paperresult}{Lemma A.6 (Decisive advantage)}
Fix ranked finite-mean classes, a reserve \(p>0\), a high cost \(c_H>0\), and a family of continuous incumbent laws.
The family has \emph{decisive advantage} if, for every law \(F\) in it,
\[
g_H(F)>c_H\quad\Longrightarrow\quad W_H(F)=1.
\]
At every law in this family with \(\DT>0\) and \(g_H>g_L\), any equilibrium that recruits a high-cost challenger after
some order has \(W_H(F)=1\). This holds at every interior prior, under all four tie rules, both at fixed presence and
with research at a cap below one. Under (H1)--(H3), decisive advantage at \(F_1\) gives the ownership and
takeover-share conclusions under outbidding in Proposition~2(ii).

For point classes \(h>\ell\ge p\) and the full family \(F_r=U[0,r]\), \(r>0\), decisive advantage holds exactly when
\[
c_H\ge\frac{h^2-p^2}{2h}.
\]
\end{paperresult}

In words, a high cost screens the incumbents against which recruitment can happen: class \(H\) earns more than
\(c_H\) only against incumbents that it beats for sure.

\begin{derivation}[proof]
\emph{Consequence.} A liquidity trader can place every order in both classes (also with research, because presence
stays below a cap below one), so every price belief satisfies \(0<\bar\mu<1\). Recruitment needs
\(B(\bar\mu)\ge c_H\) under either tie rule. Then
\begin{align*}
c_H\le B(\bar\mu)=g_H-(1-\bar\mu)(g_H-g_L)<g_H \why{\(1-\bar\mu>0\), \(g_H>g_L\)}
\end{align*}
and decisive advantage gives \(W_H(F)=1\). Under (H1), \(c_H<B_1(M)<g_{H,1}\) in the same way.

\emph{Uniform family.} \(W_H(F_r)=F_r(h)\), which is one for \(r\le h\) and \(h/r<1\) for \(r>h\). For \(r>h\),
\(g_H(F_r)=\int_p^h(y/r)\,dy=(h^2-p^2)/(2r)\), which falls strictly in \(r\) toward its supremum \((h^2-p^2)/(2h)\)
as \(r\downarrow h\). If \(c_H\ge(h^2-p^2)/(2h)\), then \(g_H>c_H\) only for \(r\le h\), where \(W_H=1\). If
\(c_H<(h^2-p^2)/(2h)\), every \(r\) in \(\bigl(h,(h^2-p^2)/(2c_H)\bigr)\) has \(g_H>c_H\) and \(W_H<1\).
\end{derivation}

\begin{chaptersummary}
Proposition~2 compares three economies that differ only in incumbent strength; each has a unique equilibrium under
all four tie rules. Against the weak incumbent no informed order covers \(k\), the price is silent, and entry is
\(\rho\). Against the strong incumbent a purchase lifts the belief to \(M(\pi)\) and recruits a high-cost
challenger, so entry rises to \(\rho+(1-\rho)(b+\pi\lambda)\). Against the very strong incumbent the price still
speaks, but no price covers \(c_H\), and entry returns to \(\rho\). High-class ownership rises exactly when
\(e_HW_H(F_1)>\rho W_H(F_0)\); decisive advantage at \(F_1\) gives \(W_H(F_1)=1\) and the clean rise.
\end{chaptersummary}
